\documentclass[11pt]{article}
\usepackage[a4paper,margin=24mm]{geometry}
\usepackage{amsmath,amssymb,amsthm}
\usepackage[T1]{fontenc}
\usepackage{lmodern}
\usepackage[hidelinks]{hyperref}
\newtheorem{theorem}{Theorem}
\newtheorem{lemma}{Lemma}
\newcommand{\Mat}{\mathcal M}
\setlength{\parindent}{0pt}
\setlength{\parskip}{5pt}
\title{A disproof of conjecture 00000000989\\\large Complete positivity excludes positive-map indecomposability}
\author{C0ldSmi1e}
\date{October 5, 2026}
\begin{document}
\begin{center}
{\LARGE A disproof of conjecture 00000000989\par}
\vspace{5pt}
{\large Complete positivity excludes positive-map indecomposability\par}
\vspace{10pt}
C0ldSmi1e \quad\textbullet\quad October 5, 2026
\end{center}

\section{Statement and mathematical conventions}
The bilingual conjecture asks for an explicit family of indecomposable
completely positive maps of dimension $d$ with Choi rank $d^2-1$.
We use the standard positive-map meaning of indecomposability, stated below.
The source does not specify the domain and codomain, dimension range, or
family index. We take a dimension-$d$ map to be a complex-linear map
$\Phi:\Mat_d(\mathbb C)\to\Mat_d(\mathbb C)$ for a positive integer $d$;
$d$ is the matrix side length. A family must have at least one member.
The obstruction applies in every positive dimension and at every Choi rank,
so no particular lower cutoff or family indexing is needed.

A matrix is positive when it is Hermitian positive semidefinite. A linear map
is positive if it preserves such matrices, and completely positive (CP) if
$\operatorname{id}_k\otimes\Phi$ is positive for every positive integer $k$.
A positive map is decomposable if it can be written
\[
 \Phi=S+\tau\circ T,
 \qquad S,T\text{ completely positive},
\]
where $\tau(X)=X^{\mathsf T}$ is ordinary matrix transpose on the output
algebra. A positive map is indecomposable if no such representation exists.
This is the convention in \cite[Section~1]{mh}. The summands may be zero;
there is no requirement of trace preservation or unitality.

This notion differs from an extreme ray of a cone, an extreme point of a
normalized convex set, or indecomposability under direct sums. The source
specifies none of those alternatives. Its parenthetical description
``extremal Choi rank'' gives no optimization class and does not change the
stated numerical equality.

For clarity, with matrix units $E_{ij}$ we use the input-first Choi convention
\[
 C_\Phi=\sum_{i,j=1}^{d}E_{ij}\otimes\Phi(E_{ij}),
 \qquad
 (C_\Phi)_{(i,k),(j,\ell)}=\Phi(E_{ij})_{k\ell}.
\]
The Choi rank is the ordinary complex matrix rank of this $d^2$-by-$d^2$
matrix, not the rank of $\Phi$ as a linear transformation. See
\cite[Section~2]{mh} for this tensor order and \cite[Section~1]{girard}
for the rank convention (with exchanged tensor factors).

\section{Complete proof}
\begin{lemma}
The zero linear map between complex matrix algebras is completely positive.
\end{lemma}
\begin{proof}
For each positive integer $k$, its ampliation sends every matrix to the zero
matrix. The zero matrix is Hermitian positive semidefinite. Thus every
ampliation preserves positive matrices, exactly as required.
\end{proof}

\newpage
\begin{theorem}
Every completely positive map between complex matrix algebras is
decomposable. Hence no such map is indecomposable in the positive-map sense.
\end{theorem}
\begin{proof}
Let $\Phi$ be completely positive. In the defining decomposition, choose
$S=\Phi$ and $T=0$. Both are completely positive, and
\[
 S+\tau\circ T=\Phi+\tau\circ0=\Phi.
\]
This is a permitted decomposition, contradicting indecomposability.
\end{proof}

In particular, for every positive integer $d$ there is no map
$\Phi:\Mat_d(\mathbb C)\to\Mat_d(\mathbb C)$ that is simultaneously
completely positive, indecomposable, and of Choi rank $d^2-1$. The rank
condition cannot restore existence: the first two conditions already
contradict each other. This conclusion does not require a rank estimate
or an unproved Choi-matrix characterization of complete positivity.

Finally, suppose a proposed family has a nonempty index set $I$, positive
dimensions $d_i$, and maps $\Phi_i$ satisfying all the stated properties.
Choose $i\in I$. The preceding theorem contradicts the properties of
$\Phi_i$. Thus no nonempty family exists, and in particular no explicit
nonempty family exists. No artificial predicate for ``explicitness'' is
needed in the argument.

\section{Formalization and verification}
The accompanying Lean project uses Lean 4.19.0 and Mathlib v4.19.0.
Its objects are actual complex-linear maps on complex matrices. Matrix
positivity has its Hermitian positive-semidefinite meaning; complete
positivity concerns all positive finite ampliation levels. Map addition,
composition, zero and transpose have their ordinary matrix meanings.
The source retains an actual Choi matrix and its complex rank in the
specialization to the conjecture, and explicitly addresses the nonempty
family conclusion.

The proof is symbolic. No numerical experiment or external computation is
needed as a mathematical premise. Verification records accompany the project
and distinguish local formal checking and independent internal review from
maintainer acceptance.

From the pinned project in \texttt{lean/}, run
\begin{verbatim}
lake build
lake env lean -DwarningAsError=true Conjecture989.lean
lake env lean -DwarningAsError=true Check.lean
\end{verbatim}
The default target builds the proof. The audit source prints the definitions,
theorem types and axiom dependencies; the accompanying records also inspect
all constants originating in the compiled proof modules, including generated
ones. The formal correspondence and exact execution identities are detailed
in the submission's verification records.

\begin{thebibliography}{9}
\bibitem{mh}
A.~M\"uller-Hermes.
\emph{Decomposable Pauli diagonal maps and Tensor Squares of Qubit Maps}.
arXiv:2006.14543v1, 2020, Sections~1--2.
\href{https://arxiv.org/abs/2006.14543v1}{arxiv.org/abs/2006.14543v1}.
\bibitem{girard}
M.~Girard et al.
\emph{On the mixed-unitary rank of quantum channels}.
Author manuscript, March 31, 2020, Section~1, equations~(1)--(3).
\href{https://cs.uwaterloo.ca/~watrous/Papers/MixedUnitaryRank.pdf}{Author-hosted manuscript}.
\end{thebibliography}
\end{document}
